\section{From experimental history to a first proposal}
\label{sec:method}

InnovationZero is a proposed deployment contract for a \emph{data-to-algorithm}
(D2A) model. Its output is a reusable learning procedure, not a prediction vector
or a set of fitted coefficients. The repository provides components with which
to investigate this contract; the inspected components do not, by themselves,
establish its empirical success. We distinguish their actual interfaces from the
architecture to be evaluated~\citep{jope,clover}.

\subsection{Three existing representations, one proposed interface}

The recorded CLOVER candidate conditions on 2,051 values: 2,048 pooled encoder
values and three query diagnostics. Its PolicyLM is a pre-normalized causal
transformer with 893,957 trainable parameters, width 128, two blocks, four
attention heads, feed-forward width 512, dropout 0.1, and a maximum program
length of 640 representation tokens. A two-layer context projection supplies
the first position; a tied output projection predicts subsequent program tokens.
The end-of-sequence state and context state jointly feed quality, feasibility,
and runtime heads. These representation tokens are not frontier-model language
tokens. The encoder may read the training targets: supervised task conditioning
is permitted, whereas access to outer-test targets is not.

Two other repository interfaces must not be conflated with CLOVER. The Candle
policy consumer accepts a frozen 2,048-value context, assembled as
$512+1{,}024+512$. Hyperion v7 separately produces a 512-dimensional problem
signature. Neither fact establishes a trained 512-dimensional
\texttt{solution\_embedding} or connects such a representation to the historical
CLOVER measurements. The bounded neural-search pipeline also seeds multiple
programs before search. It is not an implementation proof of the one-proposal
contract.

For a permitted training table $D_{\mathrm{tr}}$, with the target in its final
column, and a fixed resource card $b$, the proposed interface is
\begin{align}
 z &= E_\phi(D_{\mathrm{tr}}), &
 s &= H_\theta(z,b), \nonumber\\
 a &\sim P_\theta(\,\cdot\mid s,b), &
 M &= a.\mathrm{fit}(D_{\mathrm{tr}}). \label{eq:d2a-interface}
\end{align}
Here $z$ is the \emph{problem encoding}; $s$ is the learned
\texttt{solution\_embedding}; and $a$ is the complete learning procedure with
\texttt{fit} and \texttt{predict} interfaces. A 512-dimensional $s$ is a proposed
bottleneck to test against direct conditioning and other dimensions, not an
observed property. Source lowering and an immutable commitment to $a$ occur
\emph{before} candidate-specific execution. Fitting $M$, evaluating predictions,
and optionally exporting a fitted deployment artifact occur afterward. Those
outcomes cannot change the committed procedure or trigger a replacement.
One proposal may require many autoregressive neural calls; it is not necessarily
one forward pass.

\subsection{What the protected representation must expose}

The public contract specifies input and output shapes, permitted information,
randomness, resource bounds, supported emission targets, and refusal behavior.
It need not publish the protected grammar or its production rules. Algorithm
decomposition creates structured training objects from procedures: operations,
dependencies, fitted state, and measurable behavior. Such structure can support
learning and comparison without implying that a canonical form decides general
program equivalence. A decomposition study should test whether it improves
transfer beyond retrieval and selection from the same source archive.

The historical CLOVER decoder has finite support: its 6,640-program catalog
contains the only complete programs it can return. Consequently, for this
checkpoint, the probability of generating a program outside that catalog is
zero. This is a support limitation, not an empirical finding about all future
D2A models. A compositional successor must separately establish archive
exclusion, meaningful behavioral differences, and predictive utility.

Typed construction, successful lowering, compilation, bounded execution,
native/emitted agreement, and predictive quality are separate admission stages.
The repository supports Rust, Python, C++17, and JavaScript emission over
qualified subsets; target coverage is not uniform. A refusal is an outcome, not
an invitation to silently try another candidate. Existing parity checks use a
scaled tolerance, $10^{-9}\max(1,|y_{\mathrm{reference}}|)$ per prediction;
they do not certify bitwise equivalence of every program on every input.

\subsection{The training game and its objectives}

At the algorithm level, the decision problem is a one-step contextual bandit:
the context is a training-family problem, the action is a complete procedure,
and the consequence is its measured performance. Autoregressive construction
does not introduce intermediate environmental rewards. \emph{Algorithmic
self-play} can describe repeated training against an archive of prior
procedures, provided the archive, comparison identity, update rule, and control
are specified. A moving archive alone does not demonstrate self-improvement.

For two inner folds, let $r_j$ denote candidate RMSE and $b_j$ the RMSE of a
mean predictor fitted on the complementary rows. Historical quality is
\begin{equation}
 q(c,a)=-\frac12\sum_{j=1}^{2}
 \log\frac{\max(r_j,\varepsilon_c)}{\max(b_j,\varepsilon_c)}.
 \label{eq:inner-quality}
\end{equation}
The floor is scale-aware. The regression target clips $q$ to
$[-\log 1000,\log 1000]$; a separate rounded grid defines ties.
These inner-CV labels are training signals, not untouched generalization
estimates. The implemented native objective is
\begin{align}
 L_6={}&L_{\mathrm{CE}}+0.25L_{\mathrm{BT}}+0.10L_{\mathrm{list}}\nonumber\\
       &+0.25L_q+0.10L_{\mathrm{feas}}+0.05L_{\mathrm{rt}}.
 \label{eq:l6}
\end{align}
Every term is averaged within supported contexts, then across those contexts.
Sequence cross-entropy uniformly averages actual eligible teachers; it is not
the quality-weighted teacher loss of a separate generic implementation.
Measured generated programs supervise auxiliary heads and ranking, but do not
automatically become imitation targets in the described v1/v2 objective.
Unknown targets contribute no fabricated value. Appendix~\ref{app:methods}
gives the precise supports and reductions.

Ranking losses train the quality head, not complete-program probabilities
directly. Falling $L_6$ therefore does not establish improved deployed
proposals. Similarly, the historical best-of-eight surplus
$J_8=\mathbb{E}[\max(h_c,q_1,\ldots,q_8)-h_c]$ can improve while the first
proposal deteriorates. InnovationZero instead proposes
\begin{equation}
 J_1(\theta)=\mathbb{E}_{c,a\sim\pi_\theta(\cdot\mid E_\phi(c),b)}[u(c,a)],
 \label{eq:first-proposal-objective}
\end{equation}
with successful outcomes assigned $u=\operatorname{clip}(q,-1,1)$ and qualifying
failures assigned $-1$. This prospective utility is distinct from historical
$q_{\mathrm{opt}}$; raw error and failure counts remain necessary. Training can
compare imitation, preference updates, and an on-policy score-function estimator
of $J_1$ under matched data and compute. None is asserted here to have won.

The recorded optimizer is AdamW, initially $3\times10^{-4}$ with weight decay
0.01, bias/norm exclusions, gradient clipping at 1, and plateau halving toward
$10^{-6}$. A v2 restart retains optimizer moments and adds a 0.2-weighted exact
finite-support KL anchor to frozen v1 on selected training contexts. Restart-only
and equal-compute controls are therefore required. Dataset aliases and variants
inherit their parent-family role; encoder preprocessing and normalization are
fitted using permitted training information. Fresh-family evaluation and
independent post-commitment measurements decide whether experimental history
has improved the first proposal.
